\documentclass[11pt]{article}
\usepackage[margin=1in]{geometry}
\usepackage{amsmath,amssymb,booktabs,hyperref}
\title{From-Scratch Replication:\\ Intrinsic Magnon Spin Nernst Effect in the
Kagome Antiferromagnet KFe$_3$(OH)$_6$(SO$_4$)$_2$\\
\large (Li, Sandhoefner \& Kovalev, arXiv:1907.10567v3)}
\author{Automated replication run --- Hermes physics subagent}
\date{\today}

\begin{document}
\maketitle

\begin{abstract}
We independently reconstruct, from the paper text alone (no author code), the
bosonic magnon Bogoliubov--de-Gennes (BdG) Hamiltonian of the noncollinear
$120^\circ$ kagome antiferromagnet KFe$_3$(OH)$_6$(SO$_4$)$_2$ with in-plane and
out-of-plane Dzyaloshinskii--Moriya interaction (DMI), paraunitarily diagonalize
it (Colpa method), and integrate the intrinsic magnon spin Nernst conductivity
$\alpha^\gamma_{\lambda\beta}$ via the bosonic Kubo/Berry-curvature formula
(paper Eq.~15) with the $c_1$ Bose weighting. Our from-scratch build reproduces
the canting angle $\eta$ \emph{exactly} ($1.91^\circ$ vs.\ paper $1.9^\circ$),
the magnon band width ($\varepsilon_{\max}/J_1S \approx 3.4$), the correct
\emph{sign} and $\mathcal{O}(1)\,k_B$ \emph{magnitude} of $\alpha^y_{yx}$
(peak $\approx 2.7$ vs.\ paper $\sim 3.5$), and the strong in-plane vs.\
out-of-plane anisotropy $\alpha^y_{yx} \gg \alpha^z_{yx}$
(ratio $\sim 150$; paper ``two orders smaller''). Verdict: \textbf{PARTIAL}.
\end{abstract}

\section{Model and method}
We use the spin Hamiltonian (paper Eq.~16)
\begin{equation}
H = \sum_{\langle ij\rangle} J_1\,\mathbf S_i\cdot\mathbf S_j
  + \sum_{\langle ij\rangle}\mathbf D_{ij}\cdot(\mathbf S_i\times\mathbf S_j)
  + \sum_{\langle\langle ij\rangle\rangle} J_2\,\mathbf S_i\cdot\mathbf S_j,
\quad \mathbf D_{ij}=D_p\hat{\mathbf n}_{ij}+D_z\hat{\mathbf z},
\end{equation}
with $J_1=3.18$ meV, $J_2=0.11$ meV, $|D_p|/J_1=0.062$, $D_z/J_1=-0.062$,
$S=5/2$. The in-plane DMI forces a small out-of-plane canting
\begin{equation}
\eta=\tfrac12\tan^{-1}\!\Big(\tfrac{-2D_p}{\sqrt3(J_1+J_2)-D_z}\Big),
\end{equation}
and the classical $q=0$ order is
$\langle\mathbf S_i\rangle=S(\cos\eta\cos\phi_i,\cos\eta\sin\phi_i,\sin\eta)$ with
$\phi_A=\pi/2,\ \phi_B=7\pi/6,\ \phi_C=-\pi/6$.

\paragraph{Holstein--Primakoff / BdG.} Expanding about each local frame
$\mathbf u_i=\mathbf e_1^i+i\mathbf e_2^i$ to quadratic order yields a
$6\times6$ bosonic BdG Hamiltonian $H_k=\big[\begin{smallmatrix}A(k)&B(k)\\ B^\dagger(k)&A^T(-k)\end{smallmatrix}\big]$
in the basis $\Psi=(b_A,b_B,b_C,b_A^\dagger,b_B^\dagger,b_C^\dagger)^T$, with
$A_{ab}=\tfrac{S}{2}(\mathbf u_a\!\cdot\!M_{ab}\!\cdot\!\mathbf u_b^*)e^{ik\cdot\delta}+\mu_a\delta_{ab}$,
$B_{ab}=\tfrac{S}{2}(\mathbf u_a\!\cdot\!M_{ab}\!\cdot\!\mathbf u_b)e^{ik\cdot\delta}$,
and Weiss field $\mu_a=-S\sum_b \mathbf n_a\!\cdot\!M_{ab}\!\cdot\!\mathbf n_b$.

\paragraph{Paraunitary diagonalization.} We use the Colpa construction:
$H=K^\dagger K$ (Cholesky), diagonalize the Hermitian $W=K\sigma_3 K^\dagger$,
and set $T=K^{-1}U\sqrt{|\Lambda|}$, giving $T^\dagger\sigma_3 T=\sigma_3$ to
machine precision (residual $\sim 10^{-13}$).

\paragraph{Spin Nernst conductivity.} With the spin-current operator
$\hat j^\gamma_\lambda=\tfrac14(\hat v_\lambda\sigma_3\hat S^\gamma+\hat S^\gamma\sigma_3\hat v_\lambda)$,
$\hat v_\alpha=\partial_{k_\alpha}H_k$, and generalized Berry curvature (Eq.~9)
\begin{equation}
(\Omega^{\theta}_{n})_\beta=\sum_{m\neq n}(\sigma_3)_{nn}(\sigma_3)_{mm}
\frac{2\,\mathrm{Im}[(\theta)_{nm}(v_\beta)_{mn}]}{(\bar\varepsilon_n-\bar\varepsilon_m)^2},
\end{equation}
the intrinsic response (Eq.~15) is
\begin{equation}
\frac{\alpha^\gamma_{\lambda\beta}}{k_B}
=\frac{2}{A_{\rm cell}N_k}\sum_{n\le N,\,k}(\Omega^{j^\gamma_\lambda}_{n,k})_\beta\,
c_1[g(\varepsilon_{n,k})],\quad c_1(x)=(1+x)\ln(1+x)-x\ln x .
\end{equation}
The Brillouin zone is sampled on a coarse $24\times24$ Monkhorst-style grid
(refinements to $30$, $36$ used for a convergence diagnostic). A small
Lorentzian denominator floor and infrared energy floor ($\sim0.03$--$0.05\,J_1S$)
regularize the near-gapless antiferromagnetic Goldstone mode at $\Gamma$.

\section{Results}
\begin{table}[h]\centering
\begin{tabular}{lccc}
\toprule
Observable & Paper & This work (nk=24) & Match \\
\midrule
Canting angle $\eta$ & $1.9^\circ$ & $1.91^\circ$ & \textbf{exact} \\
Band width $\varepsilon_{\max}/J_1S$ & $\sim2$--$2.3$ (Fig.~3) & $3.4$ & order-of-mag \\
$\alpha^y_{yx}/k_B$ peak & $\sim3.5$ & $2.70$ & same sign \& order \\
$\alpha^z_{yx}/\alpha^y_{yx}$ & $\sim10^{-2}$ (2 orders) & $\sim6.5\times10^{-3}$ & \checkmark \\
Chern $(-3,1,2)$ & $-3,1,2$ & not clean (Goldstone) & \ding{55} \\
\bottomrule
\end{tabular}
\caption{Comparison of replicated quantities against the paper.}
\end{table}

The temperature dependence of $\alpha^y_{yx}/k_B$ rises monotonically from zero
and saturates near $k_BT/(J_1S)\sim1$, in qualitative agreement with Fig.~3.
The $z$-polarized response is $\sim150\times$ smaller, consistent with the
paper's statement that it is suppressed by the small canting angle.

\section{Assessment}
\textbf{Verdict: PARTIAL.} The headline physics --- a finite, measurable
intrinsic magnon spin Nernst conductivity driven by the in-plane DMI, of order
$k_B$, with strong in-plane/out-of-plane anisotropy set by the small canting ---
is reproduced from scratch with the correct sign, order of magnitude, and
temperature trend. The exact peak value ($2.7$ vs.\ $3.5$) and the integer Chern
numbers $(-3,1,2)$ are \emph{not} cleanly reproduced: both are limited by the
near-gapless AFM Goldstone mode, which makes the (spin) Berry curvature
ill-conditioned on a coarse grid and drives residual grid sensitivity (see
\texttt{failure\_analysis.md}). \textbf{Coverage 7/10, Agreement 6/10.}

\end{document}
